\documentclass{article}
\usepackage{../header}
\usepackage{datetime2}
\usepackage{xcolor}
\usepackage{colortbl}
\def\Xint#1{\mathchoice
	{\XXint\displaystyle\textstyle{#1}}%
	{\XXint\textstyle\scriptstyle{#1}}%
	{\XXint\scriptstyle\scriptscriptstyle{#1}}%
	{\XXint\scriptscriptstyle\scriptscriptstyle{#1}}%
	\!\int}
\def\XXint#1#2#3{{\setbox0=\hbox{$#1{#2#3}{\int}$ }
		\vcenter{\hbox{$#2#3$ }}\kern-0.6\wd0}}
\def\ddashint{\Xint=}
\def\dashint{\Xint-}
\usepackage[
colorinlistoftodos,
textsize=small,          % Sets the text to the smallest LaTeX font size
textwidth=0.7in        % Sets note width slightly smaller than the 1-inch margin
]{todonotes}
\definecolor{faintgray}{gray}{0.80}
\usepackage{multirow}
\rhead{tl;dr Notes on Nonlinear Wave Equations}
\om 
\title{tl;dr Notes on Nonlinear Wave Equations}
\author{\me}
\date{Fall 2026}

\begin{document}
	\maketitle
	\fullline 
	\tableofcontents
	\xappto\thefootnote{\relax}
	\footnotetext{Last updated: \today\ at \DTMcurrenttime}
	\halfline 
	\newpage 
	\setlength{\parskip}{0.55em}

\section*{Overview} 

The following is a tl;dr version of the lecture notes on nonlinear wave equations I am writing in preparation for the SBU math PhD oral exam. The purpose of these notes aren't to include every single detail/computation that is inside the main set of notes, but rather to point out the important ideas, techniques, and perspectives.

\newpage 

\section{Syllabus}

\begin{table}[h!]
	\centering 
	\begin{tabular}{|>{\centering\arraybackslash}m{4.2cm}|m{7.5cm}|}
		\hline 
		\multicolumn{2}{|c|}{\textbf{Functional and Harmonic Analysis}} \\ 
		\hline 
		\multirow{2}{*}{Harmonic Analysis} & Three-Lines Lemma \newline Riesz Interpolation Theorem   \\\hline 
	\end{tabular}
\end{table}

\begin{table}[h!]
	\centering 
	\begin{tabular}{|>{\centering\arraybackslash}m{4.2cm}|m{7.5cm}|}
		\hline 
		\multicolumn{2}{|c|}{\textbf{Nonlinear Wave Equations}} \\ 
		\hline 
		\multirow{2}{*}{Linear Wave Equation} & Cauchy Problem \newline 
		Well-Posedness \newline 
		Simple Energy Estimates \newline 
		Properties of Fundamental Solutions \newline 
		Explicit Representation Formulas \newline 
		Pointwise and Uniform Dispersion Rates ($n \geq 2$) \newline 
		Finite Speed of Propagation \newline 
		Strong Huygens Principle \\ 
		\arrayrulecolor{black!25}\cline{2-2}\arrayrulecolor{black} 
		& Hahn-Banach Theorem and Existence of Weak Solutions (Variable Coefficient Equation) \\
		\hline 
		\multirow{2}{*}{Linear Hyperbolic Systems} &  Method of Characteristics and loss of uniqueness for unbounded speed \\ \arrayrulecolor{black!25}\cline{2-2}\arrayrulecolor{black} &  Strictly/Symmetric Hyperbolic Systems \\\hline 
		Nonlinear Wave Equations & \textbf{Local} Well-Posedness \\
		\hline 
	\end{tabular}
\end{table}


\newpage
\section{Functional and Harmonic Analysis}

\begin{itemize}
	\item 
\end{itemize}

\newpage 
\section{Linear Wave Equation} 
\begin{itemize}
	\item (\textbf{Cauchy Linear Wave Problem}) Let $\phi$ be a function (or distribution) of appropriate regularity on $I \times \mathbb{R}^{n}$, and let $\Box = \partial_{tt} - \Delta$ be the d'Alembertian operator. Then the Cauchy wave problem on $I \times \mathbb{R}^{n}$ is 
	\begin{equation*}
		\left\{
		\begin{alignedat}{1}
			&\Box\phi = f,  \\
			&(\phi, \partial_{t}\phi)\upharpoonright_{\{t = 0\}} = (\phi_{0}, \phi_{1})
		\end{alignedat}
		\right., 
	\end{equation*}
	where $f$ (potentially zero) is some source term. 
	\item (\textbf{Well-Posedness of the Cauchy Problem}) A desired property for the Cauchy problem. We require the following three conditions: \vspace{-0.2cm}
	\begin{enumerate}
		\item Existence of a solution; 
		\item Uniqueness of the solution; 
		\item Continuous dependence of the solution on the initial data (small perturbation in the initial data leads to a small perturbation in the solution). 
	\end{enumerate}
	\item (\textbf{Energy of a Solution; Homogeneous}) Assuming a solution $\phi$ exists to $\Box\phi = 0$, define the energy $E(t)$ to be 
	\begin{equation*}
		E(t) = \frac{1}{2}\int_{\mathbb{R}^{n}}(\partial_{t}\phi)^{2} + (\nabla \phi)^{2}\;dx. 
	\end{equation*}
	(1) The energy is always conserved, and (2) gives uniqueness of solution (if a solution exists) by considering the difference of two solutions which must have zero energy at $t = 0$ (compare this to elliptic and parabolic equations, whose uniqueness of solutions can be given by a maximum principle). 
	\item (\textbf{Energy Estimate of a Solution; Nonhomogeneous}) Assume $\partial_{\alpha}(a^{\alpha\beta}\partial_{\beta}\phi) = F$, where $a_{\alpha\beta}$ is some metric (close enough to the standard Minkowski metric $g_{\alpha\beta}$). Let $|\partial \phi|^{2} \coloneqq (\partial_{t}\phi)^{2} + (\nabla \phi)^{2}$. Although $\norm{\partial \phi}_{C([0,T])}$ is not conserved, it is bounded above by the $\dot{H}^{1}(\mathbb{R}^{n}) \times L^{2}(\mathbb{R}^{n})$ norm of the initial data, and the $L^{2}$ norm of the source term.
	
	\item (\textbf{Fundamental Solutions}) Distributional solutions to the wave equation, 
	\begin{equation*}
		\left\{
		\begin{alignedat}{1}
			&\Box E_{+} = \delta_{0},  \\
			&\operatorname{supp}(E_{+}) \subseteq \{(t, x): 0 \leq |x| \leq t\} 
		\end{alignedat}
		\right., 
	\end{equation*}  
	I.e., a distribution $E_{+}$ solving $\Box E_{+} = \delta_{0}$ that is supported on the forward light cone based at the origin. \textbf{Key Facts:} (1) the support of the fundamental solution grows with time (i.e., solutions to the wave equation propagation), (2) \textit{if} a fundamental solution exists, then it propagates at finite speed, (3) a \textit{unique} fundamental solution exists: 
		\begin{equation*}
			E_{+}(t, x) = - \frac{\pi^{\frac{1 - n}{2}}}{2}\chi_{\{t \geq 0\}}\chi_{+}^{\frac{1 - n}{2}}\left(t^{2} - |x|^{2}\right), 
		\end{equation*} 
		where 
		\begin{equation*}
			\chi_{+}^{a}(y) \coloneqq \frac{y^{a}\chi_{\{y \geq 0\}}}{\Gamma(a + 1)}; 
		\end{equation*} 
	i.e., 
		\begin{equation*}
			E_{+}(t, x) = -\frac{\pi^{\frac{1-n}{2}}}{2}\chi_{\{t \geq 0\}}\frac{1}{\Gamma(\frac{3 - n}{2})}\frac{1}{(t^{2} - |x|^{2})^{\frac{n- 1}{2}}}\chi_{\{|x| \leq t\}},
		\end{equation*} 
	(4) can be used to derive general representation formulas for the wave equation ($\Box\phi = F$): 
	\begin{equation*}
		\phi(t, x) = -E_{+} \ast (\phi_{1}\delta_{\{t = 0\}}) - (\partial_{t}E_{+}) \ast (\phi_{0}\delta_{\{t = 0\}}) + (F\chi_{\{t\geq 0\}}) \ast E_{+},
	\end{equation*}
	(5) can be used to understand pointwise decay estimates of solutions (pointwise, for every $R > 0$ and $|x| \leq R$, $|E_{+}(t, x)| \sim O(t^{1 - n})$), (6) can be used to derive explicit representation formulas for solutions to the linear wave equation. 
	
	\item (\textbf{Explicit Representation Formulas}) Solutions exist to linear wave equation, and are unique. Suffices for the initial data $(\phi_{0}, \phi_{1}) \in H^{1}(\mathbb{R}^{n}) \times L^{2}(\mathbb{R}^{n})$. \textbf{Key Facts:} (1) if initial data are continuous and compactly supported, uniform decay rate for $n \geq 2$ is $\lesssim_{(\phi_{0}, \phi_{1})} (1 + t)^{\frac{1 - n}{2}}$, (2) also propagate in time, (3) satisfy the \textbf{Strong Huygens Principle} in odd dimensions at least 3 (solution depends only on data on spherical shells). 
\end{itemize}
	
	
  \end{document} 